\begin{table}[htbp]
\centering
\caption{Section 1115 SUD Waiver Adoption Timeline}
\label{tab:waiver_timeline}
\begin{threeparttable}
\begin{adjustbox}{max width=\textwidth}
\begin{tabular}{llrl}
\toprule
Year & Month & $N$ & States \\
\midrule
\multirow{5}{*}{2018} & July & 1 & NH \\
 & August & 1 & CA \\
 & September & 1 & WI \\
 & October & 5 & LA, MI, NC, NJ, VT \\
 & Nov--Dec & 3 & MA, NM, RI \\
\addlinespace
\multirow{3}{*}{2019} & January & 4 & AK, IL, KS, MN \\
 & February & 1 & PA \\
 & July & 1 & NE \\
\addlinespace
\multirow{3}{*}{2020} & January & 3 & DE, OR, TN \\
 & April & 1 & MT \\
 & Jul--Oct & 2 & OH, ME \\
\addlinespace
2021 & Jan, Nov & 2 & CO, NY \\
2022 & Mar, Jul, Oct & 3 & FL, IA, MS \\
2023 & June & 1 & GA \\
\addlinespace
\multicolumn{2}{l}{Early (pre-July 2018)} & 8 & AZ, IN, KY, MD, UT, VA, WA, WV \\
\bottomrule
\end{tabular}
\end{adjustbox}
\begin{tablenotes}[flushleft]
\small
\item \textit{Notes:} Treatment month $G_s$ is defined as the first month of CMS Section 1115 SUD waiver approval for each state. ``Early'' adopters whose waivers predate July 2018 are excluded from the main CS-DiD analysis (insufficient pre-treatment data in T-MSIS, which begins January 2018) but included in the ``always-treated'' robustness check. Never-treated jurisdictions (14): AL, AR, CT, DC, HI, ID, MO, ND, NV, OK, SC, SD, TX, WY. Sources: CMS State Waivers List; KFF Medicaid Waiver Tracker.
\end{tablenotes}
\end{threeparttable}
\end{table}
